\section{问题定义}
给定无向简单图 $G=(V,E)$，本文计算无序三元组 $\{u,v,w\}$ 的数量，使三个不同顶点之间的三条无向边均存在。目标总数记为
\[
T=\bigl|\{\{u,v,w\}\subseteq V:\ u,v,w\text{两两不同，且 }\{u,v\},\{u,w\},\{v,w\}\in E\}\bigr|.
\]
输入为文本边表。对于有向来源数据，固定采用“任一方向存在即保留对应无向边”的规则；去除自环，并将重复与反向重复边合并。原始顶点编号可以不连续，实验实现使用有符号64位整数读入，映射为连续32位内部下标。只出现在自环中的顶点仍保留；没有任何输入端点记录的孤立点无法从边表恢复，但不影响三角形总数。

计数和归约使用64位无符号整数，CSR偏移同样使用64位整数。当前实现未提供计数溢出检测，因此实验结果必须位于该整数范围内。本文主要比较的是相同输入、相同GPU计数功能下，基础流程与预处理优化流程完成一次任务的总时间，不以CPU串行时间替代GPU同硬件主对照。
